\chapter{Greinene}

\begin{motto}
Eit fag utan rot set likevel grein.\\ Det er berre at kvar grein ber det same fråvêret vidare og kallar det ein metode.
\end{motto}

Til no har boka følgt éi stamme: lina frå handverket gjennom modernismen og metoderørsla, fram til skjermen og ekstraksjonen. Men eit fag veks ikkje berre i tid. Det veks i breidd. Det deler seg i retningar med kvar si utdanning, kvart sitt tidsskrift, kvar sin kanon og kvar sin sjargong, og kvar med ei stille overtyding om at \emph{ho}, om ikkje dei andre, står på fast grunn. Industridesignaren held UX-aren for ein skjermpussar; UX-aren held industridesignaren for ein formgjevar av daude ting; tenestedesignaren ser ned på begge frå systemnivået; og systemtenkjaren ser ned på alle frå kartet sitt. Kvar grein trur dei andre er grunnlause og at ho sjølv ikkje er det.

Dette kapitlet stiller kvar av dei det same spørsmålet boka har stilt stammen heile vegen: kva er forma, og kva verkar på henne, så presist at to fagfolk som ikkje kjenner kvarandre ville svare likt? Svaret kjem til å falle likt fire gonger, men ikkje av di greinene er like late. Det fell likt av di dei deler eit problem ingen av dei vil sjå. Ein grein har grunn berre om ho har ein dommar ho ikkje har sauma sjølv: noko utanfor utøvaren som kan seie nei. Det finst nøyaktig eitt slikt nei i heile faget, og det dukkar opp i den greina ein minst skulle vente det.

\section{Industridesign: forma som sal, og det eine unntaket}

Den eldste greina er alt teikna i denne boka og treng berre samlast. Industridesignet voks ut av styling, og styling er sal gjort til form. Raymond Loewy laga ein leveregel av det: aldri la det gode vere godt nok, for det gode sel ikkje som det nye \autocite{loewy1951}. Strømlinja som la seg over kjøleskap og brødristarar på trettitalet, var ikkje aerodynamikk. Ho var eit teikn om fart, festa på ting som ikkje rører seg, av di teiknet selde \autocite{meikle1979}. Spør ein ærleg industridesignar i dag kva som styrer produktforma hans, og han svarar marknaden og smaken. Han har skisse og modell, brukartest og produksjon, men ingen av delane svarar på kvifor forma skal vere denne og ikkje ei anna, ut over at denne sel og denne likar oppdragsgjevaren. John Heskett skreiv den nøkterne historia til faget og fann at industridesignet aldri var ein autonom disiplin med eigne lover, men ein funksjon av produksjon og marknad, det leddet i verdikjeda som gjer ein vare ynskjeleg \autocite{heskett1980}. Det er inga skulding han la til. Det er det han fann.

Og likevel ligg det sjeldnaste i heile boka nett her, i den mest kommersielle greina. Det er alt teikna ut i kapitlet om den amerikanske stylingen: Dreyfuss la til eitt einaste omsyn som lét seg måle, menneskekroppen, og let det disiplinere forma der ho rører handa \autocite{dreyfuss1955}. Det som ber over hit, er skiljet, ikkje figuren. Joe og Josephine vart \emph{målte} og ikkje \emph{skrivne}, og difor kan dei seie nei til den som teiknar: eit handtak sit anten innanfor rekkjevidda til femte percentil eller ikkje, same kva formgjevaren føretrekkjer. Det er eit målbart brot mot noko som er som det er. Hald fast ved den skilnaden mellom det målte og det skrivne. Han kjem att i kvar einaste grein under, og kvar gong manglar han.

Dette er ikkje design som endeleg oppfører seg som ein vitskap, og kroppen er ingen mal resten av faget skulle ha nått, slik kapitlet om faget som ikkje kan seie nei held fast. Ergonomien dekkjer berre det tynne sjiktet der handa møter tingen. Alt det andre forma gjer, om foreldinga, om materialet sin lagnad, om kva tingen gjer med kulturen, ligg framleis utanfor kvart måleinstrument, like grunnlaust som før.

\section{Systemorientert design: kartet som ikkje er ein teori}

Den neste greina lovar det motsette av styling. Der industridesignet snevrar verda inn til eit produkt på ei hylle, opnar den systemorienterte designen ho så vidt som råd. Birger Sevaldson og miljøet kring han bygde ein heil retning, Systems Oriented Design, på premisset om at samtidas problem er for samanvovne til å gripast som einskildobjekt, og at designaren difor må handtere det Sevaldson kallar superkompleksitet \autocite{sevaldson2022complexity}. Verktøyet er gigamapping: enorme visualiseringar som veks utover heile veggen og knyter saman aktørar, relasjonar, materialstraumar og interesser i eitt einaste lerret \autocite{sevaldson2011giga}. Det skal seiast med ein gong at dette er kraftfullt. Eit gigakart gjer synleg det lineær tekst aldri greier: at alt heng saman med alt, at eit inngrep her forplantar seg dit, at ingen aktør står åleine. Som måte å sjå på er han framifrå, og programmet er ærleg formulert som ein praksisforsking der teikninga sjølv er ein måte å vite på \autocite{sevaldson2010discussions, morrison2010getting}.

No skal det seiast at greina sjølv hevdar meir enn dette. Sevaldson held fram at gigamappinga ikkje berre skildrar, men gjev grunnlag for å gripe inn, og at teikninga skal peike ut leverage points, dei punkta der eit inngrep verkar sterkast \autocite{sevaldson2022complexity}. Det er ein normativ og nesten prediktiv påstand, og kritikken må møte han på det. Men still kartet det spørsmålet ein slik påstand må tole, og det fell stille. Kva skal du gjere? Eit kart over hundre relasjonar gjev hundre moglege inngrep, og sjølve kartet rangerer ingen av dei. Skal ein maksimere velferd, redusere utslepp, auke rettferd, spreie makt? Aksane er teikna samstundes, og kva punkt som verkar sterkast, kjem an på kva ein vil oppnå, det kartet ikkje seier. Difor vert kvart inngrep grunngjeve med ein lekk i kartet, og leverage point-en vert vald utanfor kartet og lesen inn att. Det kjennest raust og er handlingslamt. Dette er Dreyfuss snudd på hovudet. Han hadde éin akse med ein hard målestokk; gigakartet har alle akser og ingen.

Verre er det at kartet aldri vert ferdig. Kvar node kan koplast til fleire, kvar relasjon kan brytast i finare relasjonar, og superkompleksiteten er per definisjon utan botn. Dette er den manglande stoppregelen frå kapitlet om dei vrange problema, no gjord til dyd \autocite{buchanan1992wicked}: gigamappinga feirar at kartet alltid kan bli rikare, og kallar det djupn.

Den djupaste innvendinga er likevel den enklaste, og ho er den same prøva kapitlet om dei vrange problema heldt opp mot heile faget. Eit gigakart kan ikkje vere gale. Det kan vere rikt eller fattig, men det kan ikkje takast feil, for det hevdar ingenting som kan prøvast. Det seier ikkje at gjer du X, følgjer Y; det seier at alt heng saman, og det er sant uansett. Eit kart er ei framstilling som kan vere meir eller mindre fullstendig, aldri usann, og difor er det ikkje ein teori. Du kan teikne det vakraste gigakart over eit problem og stå like rådvill framfor det, for kartet lova aldri å seie deg kva som verkar, berre at mykje gjer det. Greina som tek superkompleksiteten på største alvor, gjer fråvêret av ein modell til sjølve metoden, og merkar det ikkje, for det instrumentet som ikkje kan ta feil, seier heller aldri ifrå.

\section{Tenestedesign: ein prosess utan ein modell av det gode}

Den tredje greina flytta objektet endå ein gong, frå tingen og frå systemet til \emph{tenesta}: den koreograferte rekkja av møte mellom ein organisasjon og dei ho tener. Opphavet er ærverdig. Tidleg på åttitalet la Lynn Shostack fram \emph{service blueprint}, eit flytdiagram som teiknar tenesta som ei line av kontaktpunkt med ei synlegheitsline som skil det kunden ser frå maskineriet bak. Det var ein verkeleg oppfinning, for fyrste gong eit notasjonsspråk for noko som før var usynleg. Tre tiår seinare samla Marc Stickdorn og Jakob Schneider feltet i \emph{This Is Service Design Thinking} og gav greina prinsipp, ein verktøykasse og eit vokabular om kundereiser, touchpoints og framføring. Ezio Manzini løfta ambisjonen til samfunnsnivå og hevda at i ei verd der alle designar, kan designkompetansen styrast mot sosial innovasjon og kollektiv velferd \autocite{manzini2015}.

Korleis det går, ser ein ikkje i teorien, men i kva folk faktisk teiknar på veggen ein tysdag føremiddag. Eit leigd møterom, gule og rosa lappar i bunkar, ein fasilitator som har skrudd opp ein papirrull tvers over langveggen. Oppdraget er ei kommunal helseteneste. På rullen veks det fram eit blueprint: øvst kundereisa til «Anne, 67», ei rad blå lappar frå «får brev» til «møter på poliklinikken», under henne synlegheitslina, og under henne att det gule maskineriet av sakshandsamarar, journalsystem og ventelister. Klokka elleve er veggen vakker, og dei sju i rommet har forstått noko verkeleg. Dei ser tenesta som ei line, dei ser kvar Anne fell mellom to stolar, dei ser kvar eit brev kjem for seint. Notasjonen verkar.

Gå så nær veggen og les han. Kvar einaste lapp svarar på \emph{korleis}. Korleis steget går, korleis kontaktpunktet ser ut, korleis ein glattar friksjonen. Ikkje éin lapp svarar på om steget burde finnast. Den same veggen, med dei same blå og gule lappane og den same synlegheitslina, kunne teikna ei teneste som sorterer søkjarar ut av ei yting like presist som han teiknar ei som hjelper Anne inn i henne. Fasilitatoren har eit ord for ei reise utan friksjon, «sømlaus», og inkje ord på heile veggen for ei reise som ikkje burde gjerast sømlaus. Notasjonen er moralsk tom på same vis som perspektivteikninga er det: ho framstiller eit sjukehus og ei felle like trufast. Det folka i rommet kjenner, er ei ærleg kjensle av at noko uklart er gjort klart, for representasjonen \emph{er} betre enn inga representasjon. Det dei ikkje merkar, er at kjensla av å forstå og det å ha ein prøvbar modell er to ulike ting som lappane har smelta saman. Veggen kan ikkje ta feil. Han kan berre vere meir eller mindre fullstendig, og difor heller aldri vise seg rett.

Eit fag som har eit presist vokabular for korleis ein optimaliserer ei teneste, og inga setning for når ein ikkje skal det, vil optimalisere det som vert bestilt. Manzini sin draum om designkompetanse i velferdas teneste og marerdaumen om designkompetanse i overvakingas teneste er, mælt med vokabularet til greina sjølv, det same handverket \autocite{manzini2015}. Kundereisa er den same notasjonen anten kunden er ein innbyggjar som skal hjelpast eller ein stat som skal sortere. Greina har språket til å gjere tenesta betre. Ho manglar språket til å seie at denne tenesta ikkje burde gjerast.

\section{Interaksjonsdesign og UX: brukaren som vart oppdragsgjevarens}

Den fjerde greina hevda høgast av alle å ha funne fast grunn, og difor er fråvêret mest pinleg her når det vert synleg. Då forma flytta inn i skjermen, mangla ho den iboande lesbarheita det fysiske objektet hadde, og faget fann ein ny dommar: brukaren. Donald Norman lærte greina at ei god form gjer sine eigne bruksmoglegheiter synlege, affordans, og at den som dyttar ei dør som skulle dragast, ikkje er dum, men lurt av ei form som lyg om kva ho vil ha gjort \autocite{norman2013}. Då Norman seinare la til den emosjonelle dimensjonen, gjorde han opplevinga til endå ein akse å optimere, men la ikkje til spørsmålet om føremålet \autocite{norman2004emotional}. Jakob Nielsen gjorde brukbarheit til ein målbar storleik og kokte heuristikkane ned til ein sjekkliste éin testar kunne køyre på ein ettermiddag \autocite{nielsen1993, nielsen1990heuristic}. Alan Cooper gav greina personas, så designaren skulle byggje for ein konkret, namngjeven person og ikkje for seg sjølv \autocite{cooper1999inmates}. For fyrste gong, lova dette, kunne ei form vere god eller dårleg på ein observerbar måte og ikkje på smak. Greina fekk faktisk ein målestokk utanfor seg sjølv, og det er meir enn dei tre andre kan seie.

Sjå likevel kvar Cooper-tanken hamnar i praksis. I eit byrå med glasvegger og ståbord heng det ein laminert plakat: «Møt Martin, 34». Eit smilande arkivfoto, kjøpt for nokre kroner, av ein mann som aldri har høyrt om produktet. Under fotoet står det at Martin er prosjektleiar, har to born, er travel men teknologivand, verdset effektivitet og ynskjer å fullføre kjøp raskt og trygt. Tanken bak grepet var edel: flytt dommen ut av designaren og over i ein namngjeven person. Men Martin er ikkje funnen. Ingen har møtt han. Han vart skriven, på eit ståbord, av tre menneske som hadde eit oppdrag, og oppdraget var å auke konverteringa. Les difor «behova» hans som det dei er, ei liste ynske som alle tilfeldigvis peikar same veg som forretninga. Han ynskjer å fullføre kjøp raskt, altså nett det som tener seljaren. Kva med behova ein levande mann på 34 òg har, å bruke mindre pengar denne månaden, å sove på avgjerda, å sleppe varselet klokka tjuetre, å kjøpe \emph{mindre}? Ikkje eitt av dei står på plakaten, og dei står der ikkje av di ingen lønner eit byrå for å skrive dei inn. Dreyfuss hadde òg to figurar, men kroppen deira var målt: femte percentil rekkjer så langt og ikkje lenger, same kva formgjevaren skulle ynskje \autocite{dreyfuss1960measure}. Martin har ingen percentil. Han har ei forteljing, og forteljinga står til teneste for den som betalte plakaten.

Same plakaten møter no brukartesten, og det er der det vert tydeleg. I eit testrom med toveg-spegel prøver ein deltakar kjøpsløypa medan skjermen loggar kvart sekund, tid til fullført kjøp, talet på feilklikk, kor mange som fell av før betaling \autocite{nielsen1993}. Instrumentet måler kor \emph{glatt} vegen til kjøp er. Det måler ikkje om kjøpet burde gjerast. Set teamet i sving mot tala, og resultatet kjem av seg sjølv: median tid frå handlekorg til fullført kjøp fell frå 90 til 11 sekund, avbrotsraten halverer seg, tilfredsheita målt rett etterpå ligg høgt, og løypa går i produksjon. Tre månader seinare syner rekneskapen ein storleik ingen test fanga. Talet på angra kjøp og refusjonskrav har stige merkbart, og kundeservice rapporterer fleire som seier «eg meinte ikkje å trykkje». Ikkje noko ledd var feil etter eigne mål. Personaen sa effektivitet, testen målte effektivitet, løypa vart effektiv. Det som mangla, var ei line, på plakaten, i blueprintet, i testskjemaet, for spørsmålet om Martin \emph{burde} kjøpe så lett. Den lina finst ikkje, av di Martin vart skriven av nokon som tente på at han sa ja. Dette er Seavers felle frå ekstraksjonskapitlet, set opp eitt steg tidlegare, alt på personaplakaten \autocite{seaver2019captivating}. Greina sitt eige vokabular har ingen term for skilnaden mellom å tene ein brukar og å fange han.

Og greina vart åtvara. Ho las åtvaringa, siterte ho, og bygde vidare uendra. Lucy Suchman viste at menneskeleg handling ikkje følgjer planar slik maskinen føreset, men er situert og improvisert mot ein konkret situasjon, at brukaren ikkje er den ryddige figuren modellen treng \autocite{suchman1987}. Terry Winograd og Fernando Flores avviste at mennesket er ein informasjonshandsamar og kravde ei forståing av at vi alltid alt er kasta inn i ei verd av meining og forplikting \autocite{winograd1986}. Paul Dourish samla tråden i ein teori om kroppsleggjord samhandling, der meininga oppstår i bruk og ikkje i hovudet til ein designar som spesifiserte ho på førehand \autocite{dourish2001}. Greina kunne ha lese dette som det det var. Men kvar av desse las ho som ei klage over at modellen av brukaren var for \emph{fattig}, og ein for fattig modell kan ein gjere rikare: meir etnografi, finare metode, og så selje den finare forståinga til den same oppdragsgjevaren. Det dei tre ikkje sa rett ut, og greina difor aldri tok inn, var det hardare: at modellen ikkje berre var fattig, men \emph{interessert}, teikna av nokon, for nokon. Ein betre modell av brukaren, i tenesta til den som vil styre brukaren, er ikkje ei løysing. Det er ei oppgradering. Greina som mest stolt hevda å stå på objektiv grunn, stod heile tida på Martin, og Martin var hennar eige verk.

\section{Det same fråvêret, fire gonger}

Set greinene ved sida av kvarandre, og dei er ikkje fire forsøk på å byggje den same modellen frå ulike kantar. Dei er fire måtar å klare seg utan han på. Kvar arva tomrommet frå stammen, gav det eit nytt namn og eit nytt verktøy, og verktøyet gjorde fråvêret produktivt nok til at ingen sakna modellen. Gigakartet gjorde fråvêret av ein teori om til djupn, blueprintet gjorde fråvêret av ein etikk om til prosess, brukartesten gjorde fråvêret av ein dom om føremål om til tal, og stylinga gjorde fråvêret av eit fundament om til, rett og slett, sal. Det er den dobbelte bokføringa att, den same som opna boka og synte seg i ekstraksjonskapitlet: prosessen og vokabularet vert ført som rikdom, medan modellen som manglar, aldri kjem på balansen. Difor kjennest faget rikt nett no, med fleire greiner og fleire metodar enn nokon gong. Rikdomen er reell, men det er ein rikdom av prosess og vokabular, ikkje av kunnskap, og eit fag som hopar opp metodar utan å hope opp ein modell, ser ut til å modnast medan det berre breier seg.

Den einaste fingeren som nokon gong stakk gjennom og kjende fast grunn, var Dreyfuss sin, og berre av di han hadde det ingen av dei andre tok seg bryet med å finne: ein dommar han ikkje sauma sjølv. Kroppen vart målt. Gigakartet, blueprintet og personaen vart teikna, alle tre dommarar faget skreiv på vegner av seg sjølv, og det som er skrive av deg, kan ikkje seie nei til deg. Dette ber sviktsignaturen frå ekstraksjonskapitlet, no i tre nye instansar: eit kart som ikkje kan vere usant, ein vegg som berre kan vere meir eller mindre fullstendig, ein persona som ynskjer det oppdragsgjevaren ynskjer. Fire kompetansar til å seie ja, finare og finare, langs den eine aksen oppdraget gav, og utanom Dreyfuss sin målte kropp inkje nei å oppdrive.

\vignett

Om kvar grein er ein utveg frå fråvêret heller enn ei fylling av det, melder eit ubehageleg spørsmål seg. Korleis kan eit heilt fag, med tusenvis av kloke utøvarar, halde fire grunnlause greiner i lufta i eit hundreår utan at samanbrotet vert openbert for det sjølv? Svaret er ikkje vrangvilje. Svaret er at faget har bygt eit maskineri som held fråvêret usynleg, akkrediteringa, prisane, fagfellevurderinga, sjølve språket greinene talar, og at sjølvbedraget difor er ein stabil likevekt og ikkje ei løgn nokon fann på. Det er dette maskineriet vi no skal demontere stykke for stykke.

\laeringsmal{\begin{itemize}
\item Gjere greie for korleis kvar undergrein (industri, systemorientert, teneste, interaksjon/UX) arva fråvêret av eit fundament og gav det eit nytt namn og eit nytt verktøy.
\item Forklare kvifor Dreyfuss' ergonomi er det målbare unntaket, og kva skiljet mellom det målte og det skrivne viser om dei andre greinene.
\item Bruke «kartet er ikkje ein teori»-kritikken på systemorientert design.
\end{itemize}}

\ovingar{\begin{enumerate}
\item Vel éi av greinene. Skriv ned kva som er prosessen, kva som er vokabularet, og kva som manglar som ein delt, prøvbar modell av forma.
\item Gigamapping gjev rike bilete av kompleksitet. Kva ville det krevje å gjere eit gigakart \emph{falsifiserbart}, altså noko som kunne vere feil?
\item Samanlikn UX-faget sin «brukar» med ein klinisk pasient. Kva har det eine omgrepet som det andre manglar, og kven kan seie nei på vegner av kven?
\end{enumerate}}

\vidarelesing{\fullcite{heskett1980}; \fullcite{dreyfuss1960measure}; \fullcite{sevaldson2022complexity}; \fullcite{sevaldson2011giga}; \fullcite{manzini2015}; \fullcite{cooper1999inmates}; \fullcite{dourish2001}.}
